% ======================================================================
% THEORY SECTION — for chapters/3-Theory-and-methods.tex
% Rotor-stator interaction, blade-passing frequency, Tyler-Sofrin modes,
% resonance thresholds, and link to the thesis's design matrix.
% ======================================================================

\section{Rotor--stator interaction and pressure pulsations}
\label{sec:rsi}

In the compact \ac{RDT}--\ac{GV}--duct configuration considered in
this thesis, the rotating \ac{RDT} blades and the stationary \ac{GV}
ring are separated by a relatively small axial distance. The flow
through the \acp{GV} is therefore not steady in a strict sense even
at the nominal operating point. Two mechanisms produce unsteady
loading on the \acp{GV} and a corresponding pressure pulsation field
in the duct~\cite{brekke,arndt1990}:
\begin{enumerate}
  \item \emph{Potential interaction.} The pressure field attached to
        each rotor blade convects past the stationary vanes. It
        produces a pressure fluctuation that decays exponentially
        with axial distance but can be significant when the axial
        gap is a fraction of a blade chord.
  \item \emph{Viscous (wake) interaction.} Each rotor blade sheds a
        velocity-deficit wake. When such a wake is chopped by a
        stationary vane, the local angle of attack changes
        momentarily and produces an impulsive change in vane loading.
\end{enumerate}
Both mechanisms are periodic in the rotor reference frame at the
rotor frequency~$f_n$, and in the stator reference frame at the
blade passing frequency~$f_{\mathrm{BPF}}$ defined below. They are
the main source of deterministic unsteadiness in the present
configuration.

\subsection{Rotor frequency and blade passing frequency}
\label{sec:rsi-bpf}

With rotor rotational speed $n$~(rpm), the rotor frequency
is
\begin{equation}
    f_n \;=\; \frac{n}{60}. \label{eq:fn}
\end{equation}
A single stationary vane observes $Z_r$ rotor blades passing in each
rotor revolution. The fundamental excitation frequency seen by a
single \ac{GV} is therefore the \emph{blade passing
frequency}~\cite{brekke}
\begin{equation}
    f_{\mathrm{BPF}} \;=\; Z_r\, f_n. \label{eq:fbpf}
\end{equation}
A single rotor blade sees $Z_s$ stationary vanes per rotor revolution
and therefore experiences a periodic loading at the \emph{vane
passing frequency}
\begin{equation}
    f_{\mathrm{VPF}} \;=\; Z_s\, f_n. \label{eq:fvpf}
\end{equation}
In the special case $Z_r = Z_s$ the two frequencies coincide,
$f_{\mathrm{BPF}} = f_{\mathrm{VPF}}$. Higher harmonics
$k f_{\mathrm{BPF}}$, $k=2,3,\dots$ are typically present because
the actual pulse shape is not sinusoidal; their amplitude decreases
with $k$ but can remain non-negligible for the second and third
harmonic if the wake is
narrow~\cite{egusquiza2006,rodriguez2007}.

\subsection{Tyler--Sofrin modal decomposition}
\label{sec:rsi-tyler}

The classical analysis by Tyler and Sofrin~\cite{tyler1962} for an
annular rotor--stator stage shows that the unsteady pressure field
downstream of the stator can be written as a sum of circumferential
modes
\begin{equation}
    p'(r,\theta,z,t) \;=\; \sum_{k=1}^{\infty}\sum_{j=-\infty}^{\infty}
      \hat{p}_{k,j}(r,z)\,
      \exp\!\left[\, i\bigl(m_{k,j}\,\theta - k\,Z_r\,\Omega\,t\bigr)\,\right],
    \label{eq:tyler_sofrin_expansion}
\end{equation}
where $\Omega = 2\pi f_n$ and the circumferential mode order is
\begin{equation}
    m_{k,j} \;=\; k\,Z_r - j\,Z_s, \qquad k,j \in \mathbb{Z}.
    \label{eq:tyler_sofrin}
\end{equation}
Each combination $(k,j)$ corresponds to a pressure pattern that
rotates around the axis with phase speed $k\,Z_r\,\Omega/m_{k,j}$ and
oscillates at frequency $k f_{\mathrm{BPF}}$ at a fixed point in the
stator frame. The mode order $m_{k,j}$ determines how the pulsation
couples with the surrounding duct: an $m=0$ pattern is axisymmetric
and propagates as a plane wave, whereas patterns with $|m|\geq 1$
are spinning modes whose upstream or downstream propagation depends
on whether the frequency is above or below the cut-off frequency of
the corresponding duct mode~\cite{morse1968}.

For the fundamental BPF harmonic $k=1$ the lowest-order spatial mode
is
\begin{equation}
    m_{\min} \;=\; \min_{j\in\mathbb{Z}} \bigl| Z_r - j\,Z_s \bigr|.
    \label{eq:mmin}
\end{equation}
Three cases are of particular relevance:
\begin{itemize}
  \item If $Z_r = Z_s$, Eq.~\eqref{eq:mmin} gives $m_{\min}=0$ with
        $k=j=1$. The lowest-order mode at $f_{\mathrm{BPF}}$ is the
        \emph{axisymmetric breathing mode}: all $Z_s$ vanes are
        loaded synchronously and in phase. The pressure field is a
        plane wave that propagates along the duct without a cut-off,
        and the integrated axial force on the \ac{GV} ring adds
        coherently at $f_{\mathrm{BPF}}$.
  \item If $\gcd(Z_r,Z_s)=1$ (coprime), the minimum mode order is at
        least $|Z_r - Z_s|$ and typically $m_{\min} \geq 1$. The net
        axial force on the ring approaches zero by cancellation
        around the circumference. The remaining loading appears as a
        rotating radial force on the rotor.
  \item If $Z_r = Z_s = 7$, the baseline configuration in this
        thesis, the breathing mode $m=0$ is unavoidable in principle.
        Whether it actually dominates the pulsation field in practice
        depends on how closely the true wake pattern of the \ac{RDT}
        satisfies the $Z_r$-fold azimuthal symmetry that the
        Tyler--Sofrin argument assumes.
\end{itemize}
It is therefore an established design guideline in axial
turbomachinery to choose $Z_r$ and $Z_s$ such that low-order spatial
modes are avoided at the harmonics of concern, especially when
structural or acoustic natural frequencies lie close to
$f_{\mathrm{BPF}}$~\cite{brekke,egusquiza2006,tanaka2011}. For
reversible pump-turbines in particular,
reference~\cite{nduwayezu2021} reports that the vaneless-space
pressure pulsation amplitude depends explicitly on the runner blade
number, which supports treating $Z_s$ as a system-level design
variable and not only as a hydraulic one.

\subsection{Structural and acoustic resonance}
\label{sec:rsi-resonance}

Unsteady blade forcing becomes critical when an excitation frequency
in Eq.~\eqref{eq:tyler_sofrin_expansion} approaches a structural or
acoustic natural frequency of the system. Two coupling categories
are typically considered~\cite{egusquiza2006,tanaka2011}.

\paragraph{Structural resonance.} The \ac{GV} blade, the \ac{GV}
ring, the \ac{RDT} blades and the supporting structure each possess
a discrete set of natural frequencies $\{f_k^s\}$. A resonance
condition occurs when an excitation frequency $k\,f_{\mathrm{BPF}}$
or $k\,f_{\mathrm{VPF}}$ matches one of the $f_k^s$, \emph{and} the
spatial mode order of the excitation is compatible with the
deformation mode shape. A breathing ($m=0$) pressure pulsation will
efficiently excite axial bending or piston-like modes, while a
spinning ($m\geq 1$) pulsation will efficiently excite diametrical
modes with the matching nodal diameter count.

The \ac{GV} blade operates in water and therefore has a strong
added-mass effect. Wet natural frequencies are generally 25--40\,\%
lower than the corresponding dry values~\cite{tanaka2011}. For a
compact hydraulic \ac{GV} of the dimensions considered here, the
lowest wet bending mode is expected to lie well above 100~Hz, but
this should be verified by a separate modal analysis before a final
\ac{GV} design is committed to~\cite{valentin2015}.

\paragraph{Acoustic resonance.} In confined liquid systems, standing
pressure waves can build up if an excitation frequency coincides with
an acoustic mode of the duct. Acoustic resonance is most commonly
associated with the $m=0$ component of $f_{\mathrm{BPF}}$ because
the plane-wave mode propagates without cut-off~\cite{nicolet2007}.

\subsection{Practical magnitudes and resonance thresholds}
\label{sec:rsi-magnitudes}

Whether an unsteady loading constitutes a practical design concern
is decided by the \emph{amplitude} of the pulsation relative to the
mean loading, not by its spectral location alone. Brekke's monograph
on hydraulic turbines~\cite{brekke,brekke2010} and the IEC~60193
standard provide the following order-of-magnitude guidelines:
\begin{itemize}
  \item Normalised pressure pulsation amplitudes
        $\tilde{p}\tsub{E} = A_p / (\rho g H)$ below approximately
        $1\,\%$ at the blade passing frequency are generally
        acceptable for continuous operation.
  \item Values in the range $1$--$3\,\%$ indicate significant
        rotor-stator interaction and warrant careful design review
        and structural fatigue assessment.
  \item Values above $3$--$5\,\%$ are indicative of resonance or of
        operation close to a hydraulic instability, and usually
        require design modification.
\end{itemize}
Equivalent thresholds for the cyclic force amplitude on a single
vane are not uniformly tabulated, but reported case studies on
Francis-type machines show that fluctuation amplitudes below
$0.5\,\%$ of the mean vane load at $f_{\mathrm{BPF}}$ are typical in
safely operating units~\cite{egusquiza2006}. These reference values
are used as diagnostic thresholds in the interpretation of the
\ac{CFD} results in Chapter~\ref{ch:results}.

\subsection{Relation to the present thesis}
\label{sec:rsi-thesis}

The choice $Z_r = Z_s = 7$ of the baseline configuration places the
fundamental spatial mode at $m_{\min} = 0$ at the frequency
$f_{\mathrm{BPF}} = 7\,f_n$. The theoretical expectation is that
this configuration is \emph{a priori} unfavourable compared with
configurations using $Z_s \in \{5,6,8,9\}$, for which
$m_{\min} \geq 1$ and the net axial force on the \ac{GV} ring at
$f_{\mathrm{BPF}}$ would cancel to leading order.

The design matrix used in the simulation campaign therefore includes
$Z_s \in \{5,6,7\}$, so that this hypothesis can be tested directly
in the actual \ac{RDT}--\ac{GV} configuration and not only argued
from classical theory. Whether the breathing mode actually
manifests at a significant amplitude depends on the degree of
$Z_r$-fold azimuthal symmetry of the incoming \ac{RDT} wake and on
the spanwise matching between rotor and stator, both of which are
imperfect in practice. A diagnostic for the effective weight of the
$m=0$ mode in the computed pulsation field is introduced in
Section~\ref{sec:spec-coherence}.

% ---- Required .bib entries (add to the project bibliography) ----
% @article{tyler1962, author={Tyler, J.M. and Sofrin, T.G.},
%   title={Axial Flow Compressor Noise Studies},
%   journal={SAE Transactions}, volume={70}, year={1962}, pages={309--332}}
% @book{brekke, author={Brekke, H.},
%   title={Hydraulic Turbines: Design, Erection and Operation},
%   publisher={NTNU}, year={2015}}
% @article{brekke2010, author={Brekke, H.},
%   title={Performance and safety of hydraulic turbines},
%   journal={IOP Conf. Series}, year={2010}}
% @article{egusquiza2006,
%   author={Egusquiza, E. and Valero, C. and Liang, Q.W. and
%           Coussirat, M. and Lucino, C.},
%   title={Rotor-stator interaction in a centrifugal pump with guide vanes},
%   journal={IAHR Symp. Hydraulic Machinery and Systems}, year={2006}}
% @article{arndt1990,
%   author={Arndt, N. and Acosta, A.J. and Brennen, C.E. and Caughey, T.K.},
%   title={Experimental investigation of rotor-stator interaction in a
%          centrifugal pump with several vaned diffusers},
%   journal={ASME J. Turbomachinery}, volume={112}, number={1}, year={1990}}
% @article{rodriguez2007,
%   author={Rodriguez, C.G. and Egusquiza, E. and Santos, I.F.},
%   title={Frequencies in the vibration induced by the rotor stator
%          interaction in a centrifugal pump turbine},
%   journal={ASME J. Fluids Engineering},
%   volume={129}, number={11}, year={2007}}
% @article{tanaka2011, author={Tanaka, H.},
%   title={Vibration behaviour and dynamic stress of runners of
%          very high head reversible pump-turbines},
%   journal={Int. J. Fluid Machinery and Systems},
%   volume={4}, number={2}, year={2011}}
% @article{valentin2015,
%   author={Valentin, D. and Presas, A. and Egusquiza, E. and Valero, C.},
%   title={Influence of the added mass effect and boundary conditions
%          on the dynamic response of submerged and confined disk-like
%          structures},
%   journal={IOP Conf. Series: Earth Environ. Sci.}, year={2015}}
% @phdthesis{nicolet2007, author={Nicolet, C.},
%   title={Hydroacoustic modelling and numerical simulation of
%          unsteady operation of hydroelectric systems},
%   school={EPFL}, year={2007}, number={3751}}
% @book{morse1968, author={Morse, P.M. and Ingard, K.U.},
%   title={Theoretical Acoustics}, publisher={McGraw-Hill}, year={1968}}
% @article{nduwayezu2021, author={Nduwayezu, E. and others},
%   title={Flow instability transferability characteristics within a
%          reversible pump turbine under large guide vane opening},
%   journal={Renewable Energy}, year={2021}}
